\chapter{Lumbar puncture}

\section*{Definition and Overview}
A lumbar puncture, commonly known as a spinal tap, is an invasive medical procedure in which a needle is inserted into the subarachnoid space of the lumbar spine to collect cerebrospinal fluid (CSF) for diagnostic analysis, measure CSF opening pressure, or administer therapeutic substances. This procedure is critical in the evaluation of suspected central nervous system infections, inflammatory conditions, and subarachnoid haemorrhage when imaging is negative. It must be performed under strict aseptic conditions to prevent infection, and it is contraindicated in patients with signs of raised intracranial pressure or coagulopathy.

\section*{Epidemiology}
Lumbar puncture is a frequently performed procedure in emergency departments, neurological wards, and paediatric clinics. It is performed across all age groups, from neonates (in whom it is a key diagnostic step during a workup for neonatal sepsis) to older adults. In Australia, thousands of lumbar punctures are performed annually. The incidence of complications, particularly post-dural puncture headache (PDPH), varies widely from 10\% to 30\% depending on the size and type of needle used, with younger adult females and individuals with low body mass index having the highest susceptibility.

\section*{Aetiology and Pathophysiology}
The procedure requires navigating anatomical structures in the lower back while ensuring patient safety and minimising tissue damage.
\begin{itemize}
    \item \textbf{Anatomical targeting:} The needle is typically inserted between the L3--L4 or L4--L5 lumbar vertebrae, which lies safely below the termination of the spinal cord (conus medullaris) at the L1--L2 level in adults.
    \item \textbf{Tissue layers penetrated:} The needle passes sequentially through the skin, subcutaneous tissue, supraspinous ligament, interspinous ligament, ligamentum flavum, epidural space, dura mater, and arachnoid mater to enter the subarachnoid space.
    \item \textbf{Cerebrospinal fluid dynamics:} CSF is produced in the choroid plexuses and circulates in the subarachnoid space; measuring the opening pressure reflects intracranial pressure.
    \item \textbf{Dural puncture hole:} The physical breach of the dura mater can lead to persistent CSF leakage, causing intracranial hypotension and resultant post-dural puncture headache.
    \item \textbf{Contraindications - Space-occupying lesions:} In patients with raised intracranial pressure due to a mass, CSF removal can create a pressure gradient leading to life-threatening brain herniation.
\end{itemize}

\section*{Clinical Features}
As a procedure, clinical features relate to patient sensations during the lumbar puncture and immediate post-procedure findings.
\begin{itemize}
    \item \textbf{Procedural sensations:} A pressure-like sensation in the lower back upon needle insertion, occasionally accompanied by a transient sharp pain or electric shock sensation down one leg if a nerve root is brushed.
    \item \textbf{Post-dural puncture headache (PDPH):} A characteristic dull, throbbing headache that is orthostatic (worsens upon sitting or standing and improves when lying flat), often accompanied by nausea, vomiting, and neck stiffness.
    \item \textbf{Local pain and tenderness:} Mild soreness, bruising, or localised discomfort at the insertion site lasting for 1 to 2 days.
    \item \textbf{Normal CSF findings:} Clear, colourless fluid resembling water, with a normal opening pressure (typically 10 to 20 cm H2O in adults).
\end{itemize}

\section*{Differential Diagnosis}
It is essential to distinguish the benign complications of lumbar puncture from rare but serious adverse events.
\begin{itemize}
    \item \textbf{Post-dural puncture headache vs. meningitis:} Differentiating a benign orthostatic PDPH from post-procedure bacterial meningitis, which presents with persistent high fever, progressive neck stiffness, and photophobia.
    \item \textbf{Spinal haematoma vs. localised back pain:} Differentiating standard local soreness from an expanding epidural haematoma, which presents with progressive lower limb weakness, numbness, or bowel and bladder dysfunction.
    \item \textbf{Subarachnoid haemorrhage vs. traumatic tap:} Distinguishing blood in the CSF caused by an intracranial bleed (uniform blood across all tubes, xanthochromia) from a traumatic needle insertion (decreasing blood in successive tubes).
    \item \textbf{Cortical venous sinus thrombosis:} A rare cause of persistent headache after lumbar puncture, which is non-positional and may be accompanied by focal neurological signs.
\end{itemize}

\section*{When to Seek Medical Care}
Patients should contact their healthcare provider if they develop a persistent headache that does not improve after 24 hours of rest, or if they experience significant discomfort at the injection site.

\textbf{Call 000 (or local emergency services) for:}
\begin{itemize}
    \item A high fever, worsening neck stiffness, or severe sensitivity to light (signs of meningitis)
    \item Progressive numbness, tingling, or weakness in the legs
    \item Loss of bowel or bladder control (incontinence or urinary retention)
    \item A rapidly expanding swelling, severe bruising, or discharge at the needle insertion site
\end{itemize}

\section*{Diagnosis and Investigations}
CSF collected during the procedure is sent for laboratory testing, which is tailored to the clinical suspicion.
\begin{itemize}
    \item \textbf{Opening pressure measurement:} Determined using a manometer attached to the needle before fluid is drawn; elevated in meningitis or pseudotumor cerebri.
    \item \textbf{CSF biochemistry:} Measuring glucose and protein levels; bacterial meningitis typically shows low glucose and high protein, while viral meningitis has normal glucose.
    \item \textbf{CSF microscopy and cytology:} Cell count and differential to detect pleocytosis (increased white blood cells) and abnormal cells.
    \item \textbf{Microbiology and PCR:} Gram stain, bacterial culture, and polymerase chain reaction (PCR) assays for viral (e.g. HSV, enterovirus) and atypical pathogens.
    \item \textbf{Immunology (Oligoclonal bands):} Electrophoresis of CSF and serum to detect oligoclonal bands, crucial for diagnosing multiple sclerosis.
\end{itemize}

\section*{Management}
Management involves preparing the patient, performing the procedure under sterile technique, and managing any post-procedural complications.
\begin{itemize}
    \item \textbf{Aseptic technique and positioning:} Placing the patient in a lateral recumbent (fetal) position or sitting position, followed by skin disinfection and local anaesthesia.
    \item \textbf{Needle selection:} Using small-gauge (22G or smaller) atraumatic (pencil-point) needles to significantly reduce the incidence of PDPH.
    \item \textbf{Post-procedure rest:} Instructing the patient to lie flat for a brief period (e.g. 1 to 2 hours) and maintain adequate hydration.
    \item \textbf{Conservative PDPH treatment:} Bed rest, oral analgesics (paracetamol, NSAIDs), and caffeine intake (which causes cerebral vasoconstriction).
    \item \textbf{Epidural blood patch:} For refractory PDPH, injecting a small volume of the patient's own venous blood into the epidural space at the puncture site to seal the dural leak.
\end{itemize}

\section*{Complications}
\begin{itemize}
    \item Post-dural puncture headache (PDPH), the most common complication.
    \item Local pain, tenderness, and bruising at the needle insertion site.
    \item Infection, including epidural abscess, discitis, or bacterial meningitis.
    \item Spinal epidural or subdural haematoma, potentially leading to spinal cord compression.
    \item Cerebral herniation, a rare but catastrophic event in patients with pre-existing raised intracranial pressure.
\end{itemize}

\section*{Prognosis and Outcomes}
The prognosis after a lumbar puncture is highly favourable. Most patients experience only mild, transient discomfort that resolves within a few days. Post-dural puncture headaches, when they occur, usually resolve spontaneously within a week or respond immediately to an epidural blood patch. The diagnostic yield of CSF analysis is extremely high, providing critical, life-saving information that guides treatment for infections, autoimmune disorders, and neurological conditions.

\section*{Prevention and Self-Care}
\begin{itemize}
    \item \textbf{Atraumatic needles:} Request the use of atraumatic (pencil-point) needles, which minimise dural tearing.
    \item \textbf{Post-procedure rest:} Remain lying flat for the recommended time immediately after the procedure.
    \item \textbf{Hydration:} Drink plenty of fluids (including water, tea, or coffee) after the procedure to help replenish CSF.
    \item \textbf{Avoid heavy lifting:} Restrict strenuous activity, heavy lifting, or vigorous exercise for 24 to 48 hours.
    \item \textbf{Monitor the puncture site:} Keep the dressing clean and dry, and inspect the site daily for redness, swelling, or leaking fluid.
\end{itemize}

\section*{Sources and Further Reading}
The following reputable organisations provide further information on lumbar puncture:
\begin{itemize}
    \item Healthdirect Australia. \textit{Lumbar puncture overview.} https://www.healthdirect.gov.au/
    \item Mayo Clinic. \textit{Spinal tap (lumbar puncture) procedure details.} https://www.mayoclinic.org/
    \item NHS. \textit{Lumbar puncture: What happens and recovery.} https://www.nhs.uk/
    \item Queensland Health. \textit{Patient guide to lumbar puncture.} https://www.health.qld.gov.au/
    \item MSD Manual. \textit{How to do a lumbar puncture.} https://www.msdmanuals.com/
\end{itemize}
